% chapters/00-abstract.tex -- Реферат (українською та іноземною мовою).
%
% Вимоги: ДСТУ 3008:2015, п. 4.3; для магістерської дисертації --
% додатково розділ "Реферат" Рекомендацій КПІ (до 500 слів, певна
% послідовність пунктів). Замініть приклад нижче своїм текстом,
% зберігаючи структуру.
\kpithesisUnnumberedChapter{Реферат}

Пояснювальна записка до \workType: \pageref{LastPage}~с.,
\arabic{figure}~рис., \arabic{table}~табл., 2~додатки, 00~джерел.
% ^^^ \pageref{LastPage} працює, лише якщо в кінці документа є мітка
% \label{LastPage} (вона вже додана в main.tex). Кількість джерел
% після підключення бібліографії (Фаза 5) можна порахувати вручну.

\textbf{Об'єкт дослідження} -- \dots

\textbf{Мета роботи} -- \dots

\textbf{Методи дослідження} -- \dots

\textbf{Наукова новизна} (для магістерської дисертації) -- \dots

\textbf{Практичне значення} -- \dots

\bigskip
\noindent\textbf{Ключові слова:} ключове слово 1, ключове слово 2,
ключове слово 3.

\bigskip
\begin{english}
  \textbf{Abstract}

  Explanatory note to the \dots: \pageref{LastPage}~p., \arabic{figure}~fig.,
  \arabic{table}~tables, 2~appendices, 00~references.

  \textbf{Object of research} -- \dots

  \textbf{Purpose} -- \dots

  \textbf{Methods} -- \dots

  \textbf{Scientific novelty} -- \dots

  \textbf{Practical value} -- \dots

  \bigskip
  \noindent\textbf{Keywords:} keyword one, keyword two, keyword three.
\end{english}
